\documentclass{article}
\usepackage{graphicx}
\usepackage{amsmath}
\usepackage{amssymb}
\usepackage{listings}
\usepackage[margin=1in]{geometry}
\usepackage{xcolor}
\usepackage{tikz}
\usetikzlibrary{shapes, arrows, positioning, calc}
\usepackage{tcolorbox}

% Define colors for syntax highlighting
\definecolor{codegreen}{rgb}{0,0.6,0}
\definecolor{codegray}{rgb}{0.5,0.5,0.5}
\definecolor{codepurple}{rgb}{0.58,0,0.82}
\definecolor{backcolour}{rgb}{0.95,0.95,0.92}

% Python style for highlighting
\lstdefinestyle{pythonstyle}{
    backgroundcolor=\color{backcolour},   
    commentstyle=\color{codegreen},
    keywordstyle=\color{magenta},
    numberstyle=\tiny\color{codegray},
    stringstyle=\color{codepurple},
    basicstyle=\ttfamily\footnotesize,
    breakatwhitespace=false,         
    breaklines=true,                 
    captionpos=b,                    
    keepspaces=true,                 
    numbers=left,                    
    numbersep=5pt,                  
    showspaces=false,                
    showstringspaces=false,
    showtabs=false,                  
    tabsize=2
}
\lstset{style=pythonstyle}

\title{Lecture 3: Multi-Layer Perceptrons and Learning XOR}
\author{Deep Learning Class}
\date{\today}

\begin{document}
\maketitle

\section{Introduction: Beyond Linear Models}

In our previous lectures, we discussed linear models for regression and classification. While powerful and interpretable, linear models have a fundamental limitation: they can only learn linear decision boundaries. Today, we'll explore how multi-layer perceptrons (MLPs) overcome this limitation, using the classic XOR problem as our motivating example.

\begin{tcolorbox}[colback=yellow!5!white, colframe=orange!50!black, title=Key Insight]
The XOR problem historically demonstrated that single-layer perceptrons cannot solve certain simple problems. This limitation led to the ``AI winter'' until the development of multi-layer networks with backpropagation.
\end{tcolorbox}

\section{The XOR Problem: A Fundamental Challenge}

\subsection{What is XOR?}

The XOR (exclusive OR) function is a simple boolean operation where the output is true only when exactly one input is true:

\begin{center}
\begin{tabular}{cc|c}
$x_1$ & $x_2$ & XOR \\
\hline
0 & 0 & 0 \\
0 & 1 & 1 \\
1 & 0 & 1 \\
1 & 1 & 0 \\
\end{tabular}
\end{center}

\subsection{Why Linear Models Fail}

A linear classifier attempts to separate classes using a hyperplane:
\[f(x) = w_1x_1 + w_2x_2 + b\]

For XOR, we need:
\begin{align}
f(0,0) &< 0 \quad \Rightarrow \quad b < 0 \\
f(1,1) &< 0 \quad \Rightarrow \quad w_1 + w_2 + b < 0 \\
f(0,1) &> 0 \quad \Rightarrow \quad w_2 + b > 0 \\
f(1,0) &> 0 \quad \Rightarrow \quad w_1 + b > 0
\end{align}

From constraints (3) and (4): $w_1 > -b$ and $w_2 > -b$

Adding these: $w_1 + w_2 > -2b$

Since $b < 0$ (from constraint 1), we have $-2b > 0$, thus $w_1 + w_2 > 0$.

But this contradicts constraint (2)! Therefore, \textbf{no linear function can learn XOR}.

\section{Multi-Layer Perceptrons: The Solution}

\subsection{Architecture}

An MLP consists of:
\begin{itemize}
    \item \textbf{Input layer}: Receives the input features
    \item \textbf{Hidden layers}: Transform the input through learned representations
    \item \textbf{Output layer}: Produces the final prediction
\end{itemize}

For a single hidden layer MLP:
\begin{align}
\mathbf{h} &= \sigma(\mathbf{X}\mathbf{W}^{(1)} + \mathbf{b}^{(1)}) \quad \text{(hidden layer)} \\
\mathbf{o} &= \mathbf{h}\mathbf{W}^{(2)} + \mathbf{b}^{(2)} \quad \text{(output layer)}
\end{align}

where $\sigma$ is a non-linear activation function.

\subsection{The Critical Role of Non-linearity}

\begin{tcolorbox}[colback=red!5!white, colframe=red!50!black, title=Why Non-linearity is Essential]
Without non-linear activation functions, multiple linear layers collapse into a single linear transformation:
\begin{align}
\mathbf{o} &= (\mathbf{X}\mathbf{W}^{(1)} + \mathbf{b}^{(1)})\mathbf{W}^{(2)} + \mathbf{b}^{(2)} \\
&= \mathbf{X}\mathbf{W}^{(1)}\mathbf{W}^{(2)} + \mathbf{b}^{(1)}\mathbf{W}^{(2)} + \mathbf{b}^{(2)} \\
&= \mathbf{X}\mathbf{W} + \mathbf{b} \quad \text{(still linear!)}
\end{align}
\end{tcolorbox}

\section{Activation Functions}

\subsection{Sigmoid Function}

The sigmoid was historically popular:
\[\sigma(x) = \frac{1}{1 + e^{-x}}\]

\textbf{Problems with Sigmoid:}
\begin{itemize}
    \item \textbf{Vanishing gradients}: $\sigma'(x) = \sigma(x)(1-\sigma(x))$ becomes very small for large $|x|$
    \item \textbf{Not zero-centered}: Can slow down gradient descent
    \item \textbf{Computationally expensive}: Requires exponential calculation
\end{itemize}

\subsection{ReLU (Rectified Linear Unit)}

The modern standard activation function:
\[\text{ReLU}(x) = \max(0, x)\]

\textbf{Advantages of ReLU:}
\begin{itemize}
    \item \textbf{No vanishing gradient} for positive inputs
    \item \textbf{Computationally efficient}: Simple max operation
    \item \textbf{Sparse activation}: Many neurons output exactly zero
    \item \textbf{Empirically successful}: Works well in practice
\end{itemize}

\section{Solving XOR with an MLP}

\subsection{Intuitive Solution}

We can think of solving XOR by learning intermediate features:
\begin{itemize}
    \item Hidden neuron 1: Detects ``$x_1 = 1$ AND $x_2 = 0$''
    \item Hidden neuron 2: Detects ``$x_1 = 0$ AND $x_2 = 1$''
    \item Output: Combines these with OR operation
\end{itemize}

\subsection{Mathematical Solution}

One possible weight configuration:
\begin{align}
\mathbf{W}^{(1)} &= \begin{pmatrix} 1 & -1 \\ -1 & 1 \end{pmatrix}, \quad \mathbf{b}^{(1)} = \begin{pmatrix} -0.5 \\ -0.5 \end{pmatrix} \\
\mathbf{W}^{(2)} &= \begin{pmatrix} 1 \\ 1 \end{pmatrix}, \quad b^{(2)} = -0.5
\end{align}

This creates decision boundaries that correctly classify all XOR patterns.

\section{Training Deep Networks: Practical Considerations}

\subsection{Capacity Control}

To prevent overfitting, we use several techniques:

\begin{itemize}
    \item \textbf{Dropout}: Randomly ``drop'' neurons during training
    \begin{lstlisting}[language=Python]
def dropout(X, p=0.5):
    mask = np.random.rand(*X.shape) > p
    return X * mask / (1 - p)  # Scale to maintain expected value
    \end{lstlisting}
    
    \item \textbf{Weight Decay}: Add L2 penalty to loss function
    \item \textbf{Batch Normalization}: Normalize activations within mini-batches
\end{itemize}

\subsection{Residual Connections}

For very deep networks, residual connections help gradient flow:
\[\mathbf{y} = \mathbf{x} + f(\mathbf{x})\]

This allows the network to learn the ``residual'' $f(\mathbf{x})$ rather than the full transformation.

\section{Key Takeaways}

\begin{enumerate}
    \item \textbf{Linear models are limited}: Cannot solve problems like XOR that aren't linearly separable
    \item \textbf{MLPs overcome this limitation}: By stacking layers with non-linear activations
    \item \textbf{Non-linearity is crucial}: Without it, deep networks collapse to linear models
    \item \textbf{ReLU is preferred}: Avoids vanishing gradients and is computationally efficient
    \item \textbf{Capacity control is important}: Use dropout, weight decay, and batch norm
    \item \textbf{Deep learning = representation learning}: Hidden layers automatically discover useful features
\end{enumerate}

\begin{tcolorbox}[colback=green!5!white, colframe=green!50!black, title=Historical Note]
The XOR problem was so significant that when Minsky and Papert (1969) proved perceptrons couldn't solve it, AI research funding dried up for years. The development of backpropagation for training MLPs in the 1980s revived the field and led to today's deep learning revolution.
\end{tcolorbox}

\section{Exercises}

\begin{enumerate}
    \item Prove that a 2-2-1 MLP (2 inputs, 2 hidden neurons, 1 output) is the minimum architecture needed to solve XOR.
    \item Implement gradient descent training for an MLP on XOR from scratch.
    \item Experiment with different activation functions (sigmoid, tanh, ReLU) and compare convergence speed.
    \item Visualize the learned hidden layer representations - what features do they encode?
    \item Try solving other non-linearly separable problems (e.g., circle inside a square).
\end{enumerate}

\end{document}